\documentclass[11pt,a4paper]{article}
% Common preamble for the hint sheets (\input from each file)
\usepackage[utf8]{inputenc}
\usepackage[T1]{fontenc}
\usepackage{lmodern}
\usepackage{amsmath,amssymb,amsthm,mathtools}
\usepackage[a4paper,margin=2.3cm]{geometry}
\usepackage{enumitem}
\usepackage{xcolor}
\usepackage[hidelinks]{hyperref}
\renewcommand{\leq}{\leqslant}
\renewcommand{\le}{\leqslant}
\renewcommand{\geq}{\geqslant}
\renewcommand{\ge}{\geqslant}
\setlength{\parindent}{0pt}
\setlength{\parskip}{0.35em}
% \hintstitle{Sheet title}{Class work / Homework}
\newcommand{\hintstitle}[2]{\begin{center}{\LARGE\bfseries Hints: #1\par}\vspace{0.5em}{\large #2\par}\end{center}
\noindent\textit{The hints to each problem are cumulative. Read one, try again, and open the next only if you are still stuck.}\par\medskip\hrule\medskip}
\newcommand{\problem}[1]{\par\bigskip{\large\bfseries Problem #1}\par\nopagebreak\smallskip}
\newcommand{\subpart}[1]{\par\smallskip\textbf{(#1)}\par\nopagebreak}
\newcommand{\hint}[1]{\par\smallskip\noindent\textcolor{black!70}{\textbf{Hint #1.}}\enspace}
\newcommand{\note}[1]{\par\smallskip\noindent\textit{Note.} #1\par}
\providecommand{\lt}{<}
\providecommand{\gt}{>}
\providecommand{\tightlist}{\setlength{\itemsep}{0pt}\setlength{\parskip}{0pt}}
\emergencystretch=3em

\begin{document}
\hintstitle{Linear algebra -- 8. Block matrices and determinants of special shape}{Class work}

\problem{1}
\hint{1}
Schur formula: \(\begin{pmatrix}A&B\\C&D\end{pmatrix}=\begin{pmatrix}E&0\\CA^{-1}&E\end{pmatrix}\begin{pmatrix}A&B\\0&D-CA^{-1}B\end{pmatrix}\). Compute the determinant of \(\begin{pmatrix}E&-B\\A&E\end{pmatrix}\) using either diagonal block as the pivot: you get \(\det(E+BA)\) and \(\det(E+AB)\).
\hint{2}
Commuting \(C,D\): \(\begin{pmatrix}A&B\\C&D\end{pmatrix}\begin{pmatrix}D&0\\-C&E\end{pmatrix}=\begin{pmatrix}AD-BC&B\\0&D\end{pmatrix}\). Cancel \(\det D\) if \(D\) is invertible; otherwise replace \(D\) by \(D+\varepsilon E\), which still commutes with \(C\), and let \(\varepsilon\to0\).
\hint{3}
Counterexample (\(2\times2\) blocks): \(A=E\), \(B=E_{21}\), \(C=E_{12}\), \(D=E_{11}\). Here \(\det M=\det(D-CB)=0\) but \(\det(AD-BC)=-1\).

\problem{2}
\hint{1}
Note the clash of notation: here the block \(E\) of \(M^{-1}\) is not the identity. Write \(I\) for the identity. Multiply \(M\) by the block-triangular matrix \(\begin{pmatrix}I&F\\0&H\end{pmatrix}\), built from the right block column of \(M^{-1}\).
\hint{2}
Since \(MM^{-1}=I\), we have \(AF+BH=0\) and \(CF+DH=I\). So the product is \(\begin{pmatrix}A&0\\C&I\end{pmatrix}\). Take determinants.

\problem{3}
\hint{1}
Let \(X=ABCD\). Then \(X=A(BCD)=AA^T\) is symmetric positive semidefinite. Find a polynomial identity for \(X\).
\hint{2}
Transpose \(A^T=BCD\) and substitute the other three relations: \(A=D^TC^TB^T=(ABC)(DAB)(CDA)\), i.e.\ \(X^2A=A\).
\hint{3}
Hence \(X^3=X^2AA^T=AA^T=X\). A positive semidefinite matrix with \(X^3=X\) has eigenvalues in \(\{0,1\}\) and is diagonalisable, so \(X^2=X\).

\problem{4}
\hint{1}
Since \(n\) is odd, \(i\mapsto2i\) is a bijection of \(\mathbb Z/n\). Relabel the indices by it: \(A\) becomes the circulant \(2E+P+P^{-1}\), where \(P\) is the cyclic shift.
\hint{2}
\(2E+P+P^{-1}=P^{-1}(E+P)^2\), and \(\det(E+P)=\prod_k(1+\varepsilon^k)=1-(-1)^n=2\). Answer: \(\det A=4\).

\problem{5}
\hint{1}
\(C_n\) is the circulant with symbol \(f(x)=1+x+x^2=(x-\omega)(x-\omega^2)\), \(\omega=e^{2\pi i/3}\). So \(\det C_n=\prod_{k}f(\varepsilon^k)\). If \(3\mid n\), some \(\varepsilon^k=\omega\).
\hint{2}
\(\prod_k(\varepsilon^k-z)=(-1)^n(z^n-1)\) because \(\prod_k(z-\varepsilon^k)=z^n-1\). Hence \(\det C_n=(\omega^n-1)(\omega^{2n}-1)=|\omega^n-1|^2\), which is \(3\) when \(3\nmid n\).
\end{document}
